Considereu el sistema d'equacions lineals següent, que està format per tres plans a l'espai
i depèn del paràmetre real $m$:
\[
\left.\begin{aligned}
x+my+z&=4\\
x+3y+z&=5\\
mx+y+z&=4
\end{aligned}\ \right\}.
\]

\begin{apartats}

\apartat{1}
Discutiu el sistema per als diferents valors del paràmetre $m$.

\begin{solucio}
La matriu de coeficients i la matriu ampliada del sistema són
\[
A=\begin{pmatrix}1&m&1\\1&3&1\\m&1&1\end{pmatrix}
\quad\text{i}\quad
\bar A=\left(\begin{array}{ccc|c}1&m&1&4\\1&3&1&5\\m&1&1&4\end{array}\right).
\]
Calculant el determinant de la matriu $A$, tenim
$\begin{vmatrix}1&m&1\\1&3&1\\m&1&1\end{vmatrix}=3+m^2+1-3m-1-m=m^2-4m+3$.
Per obtenir els valors crítics fem
$|A|=0\iff m^2-4m+3=0\iff m=\frac{4\pm\sqrt{16-12}}{2}=3,\ 1$.\\
En el cas $m\neq1,3$, tenim $|A|\neq0$ i $\operatorname{rang}(A)=\operatorname{rang}(\bar A)=3$;
com que hi ha tres incògnites, el sistema és \textbf{compatible determinat}.\\
En el cas $m=3$, tenim $|A|=0$ i $\begin{vmatrix}1&3\\3&1\end{vmatrix}\neq0$, per tant
$\operatorname{rang}(A)=2$. D'altra banda,
$\begin{vmatrix}1&3&4\\1&3&5\\3&1&4\end{vmatrix}=\begin{vmatrix}1&3&4\\0&0&1\\3&1&4\end{vmatrix}
=-\begin{vmatrix}1&3\\3&1\end{vmatrix}=8\neq0$,
i deduïm que $\operatorname{rang}(\bar A)=3$ i que el sistema és \textbf{incompatible}.\\
Finalment, en el cas $m=1$, tenim $|A|=0$, $\begin{vmatrix}1&1\\1&3\end{vmatrix}\neq0$,
$\operatorname{rang}(A)=2$, i també $\operatorname{rang}(\bar A)=2$ per tenir la primera i
la tercera files iguals. El sistema és \textbf{compatible indeterminat} amb un grau de
llibertat.

\textit{Pauta oficial:} 0,25 per calcular el determinant i trobar els dos valors crítics, i
0,25 per la discussió de cadascun dels tres casos.
\end{solucio}

\apartat{1}
Interpreteu geomètricament aquest sistema per a tots els valors del paràmetre $m$ i
resoleu-lo, si és possible, per al cas $m=1$.

\begin{solucio}
Es tracta de tres plans a l'espai, que estan en diverses posicions relatives segons el valor
del paràmetre. Si $m\neq1,3$, el sistema és compatible determinat: els tres plans es tallen
en un \textbf{únic punt} comú. Si $m=3$, és incompatible: no tenen cap punt en comú, cosa que
es veu directament perquè el primer pla, $x+3y+z=4$, i el segon, $x+3y+z=5$, són
\textbf{paral·lels i no coincidents}. Si $m=1$, el sistema és compatible indeterminat amb un
grau de llibertat: els tres plans es tallen en una \textbf{recta} comuna, que trobem resolent
el sistema:
\[
\left.\begin{aligned}x+y+z&=4\\x+3y+z&=5\\x+y+z&=4\end{aligned}\right\}
\overset{F_1=F_3}{\iff}
\left.\begin{aligned}x+y+z&=4\\x+3y+z&=5\end{aligned}\right\}
\overset{z=\lambda}{\implies}
\left.\begin{aligned}x+y&=4-\lambda\\x+3y&=5-\lambda\end{aligned}\right\}
\]
\[
x=\frac{\begin{vmatrix}4-\lambda&1\\5-\lambda&3\end{vmatrix}}{\begin{vmatrix}1&1\\1&3\end{vmatrix}}
 =\frac{12-3\lambda-5+\lambda}{2}=\frac{7-2\lambda}{2},\qquad
y=\frac{\begin{vmatrix}1&4-\lambda\\1&5-\lambda\end{vmatrix}}{\begin{vmatrix}1&1\\1&3\end{vmatrix}}
 =\frac{5-\lambda-4+\lambda}{2}=\frac12.
\]
Les solucions són tots els punts $\left(\frac{7-2\lambda}{2},\frac12,\lambda\right)$ amb
$\lambda\in\mathbb{R}$: la intersecció dels tres plans és la recta
$(x,y,z)=\left(\frac72,\frac12,0\right)+\lambda(-1,0,1)$.

\textit{Pauta oficial:} 0,25 per cadascuna de les tres interpretacions geomètriques i 0,25
per la resolució del sistema en el cas $m=1$.
\end{solucio}

\apartat{0,5}
Per a $m=1$, és possible afegir una quarta equació de manera que el sistema resultant sigui
compatible determinat i tingui com a solució $(x,y,z)=\left(3,\frac12,\frac12\right)$?
Raoneu la resposta.

\begin{solucio}
Per a $m=1$, el sistema té infinites solucions, una de les quals és
$(x,y,z)=\left(3,\frac12,\frac12\right)$: precisament la que correspon a $\lambda=\frac12$.
Per tant, qualsevol nova equació que també tingui aquest punt com a solució (per exemple,
$x=3$) i que doni una matriu de coeficients de rang 3, convertirà el sistema en compatible
determinat amb el punt desitjat com a única solució. L'equació proposada ho compleix:
\[
\operatorname{rang}\begin{pmatrix}1&1&1\\1&3&1\\1&1&1\\1&0&0\end{pmatrix}
=\operatorname{rang}\begin{pmatrix}1&1&1\\1&3&1\\1&0&0\end{pmatrix}=3.
\]
\textit{Pauta oficial:} 0,5 per donar l'equació demanada i raonar que és vàlida.
\end{solucio}

\end{apartats}
